\documentclass[10pt,a4paper]{amsart}
\usepackage[margin=19mm]{geometry}
\usepackage{amsmath,amssymb,mathtools}
\usepackage[scr=rsfs,cal=euler]{mathalfa}
\usepackage{xfrac}
\usepackage[unicode,hidelinks,bookmarks=false]{hyperref}
\newcommand{\ie}{i.e.\ }
\newcommand{\cf}{cf.\ }
\newcommand{\resp}{respectively\ }
\newcommand{\Subsection}{\subsection}
\providecommand{\nobreakdash}{\nobreak\hbox{-}}
\setlength{\emergencystretch}{2em}
\multlinegap0pt
\usepackage[shortlabels]{enumitem}
\usepackage{mathtools}
\newtheorem{assumption}{Assumption}
\newtheorem{theorem}{Theorem}
\newtheorem{lemma}{Lemma}
\newcommand{\EE}{\mathbb{E}}
\newcommand{\RR}{\mathbb{R}}
\newcommand{\Ystar}{\mathcal{Y}^*}
\newcommand{\afrak}{\mathfrak{a}}
\newcommand{\bfrak}{\mathfrak{b}}
\newcommand{\xstar}{x^*}
\newcommand{\ystar}{y^*}
\title{Non-Expansive Mappings in Two-Time-Scale Stochastic Approximation: Finite-Time Analysis}
\author{Siddharth Chandak}
\date{Source: arXiv:2501.10806v4 (CC BY 4.0)}
\begin{document}
\maketitle

\noindent\emph{Source and licence.} Non-Expansive Mappings in Two-Time-Scale Stochastic Approximation: Finite-Time Analysis. Version arXiv:2501.10806v4 (CC BY 4.0), distributed by arXiv under the Creative Commons Attribution 4.0 International licence, http://creativecommons.org/licenses/by/4.0/, as recorded in the arXiv licence metadata for this exact version and shown on the abstract page https://arxiv.org/abs/2501.10806v4. Full licence notice: \texttt{licenses/arxiv-2501.10806-CC-BY-4.0.txt}.\par\smallskip

\noindent\textit{Editorial scope.} Counted unit: the article's theorem thm:main (source0.tex lines 255--262). The article's complete original proof of it is delivered with it (source0.tex lines 893--918, one \texttt{proof} environment of 26 lines, which the article places away from the statement). No mathematics is written, repaired or re-derived: the statement and the proof are the article's own lines, in the article's own notation, and no retained line is substituted or re-flowed.  Retained uncounted as the source-local context the proof needs: lines 201--206; lines 210--214; lines 247--251; lines 364--372; lines 823--827; lines 851--853.  Nothing was omitted from any retained span, so the package carries no omission record.  No retained line contains a citation, so the wrapper carries no bibliography and no bibliography entry.  No retained line contains a figure environment, an image-inclusion command or drawing code, so no image is delivered and the package declares no asset dependency.  Editorial formatting only, in addition: an AMS \texttt{amsart} preamble that restates the article's own declarations and macros used by the retained lines, namely the environments \texttt{assumption}, \texttt{theorem}, \texttt{lemma} on separate counters, together with the standard packages those lines need.  The retained environments print in this document as Assumption 1, Theorem 1, Lemma 1 in order of appearance, whereas the article prints the same items with the numbering of its own full document.  The complete native source of the article as distributed in the arXiv e-print for this exact version is bundled unchanged as \texttt{sources/171773/source0.tex}.
\begin{equation}\label{iter-main}
    \begin{aligned}
    &x_{k+1}=x_k+\alpha_k(f(x_k,y_k)-x_k+M_{k+1})\\
    &y_{k+1}=y_k+\beta_k(g(x_k,y_k)-y_k+M'_{k+1}).
\end{aligned}
\end{equation}

\begin{assumption}\label{f-contrac}
The function $f:\RR^{d_1}\times \RR^{d_2}\mapsto \RR^{d_1}$ is contractive in $x$ for each $y\in\RR^{d_2}$. Specifically, there exists a constant $0\leq \mu<1$, independent of $y$, such that
   $$\|f(x_1,y)-f(x_2,y)\|\leq \mu \|x_1-x_2\|,$$
for all $x_1,x_2\in\RR^{d_1}$ and $y\in\RR^{d_2}$.
\end{assumption}

\begin{assumption}\label{assu:stepsize}
The step sizes $\alpha_k$ and $\beta_k$ are of the form $$\alpha_k=\frac{\alpha}{(k+K_1)^{\afrak}}\;\;\text{and}\;\;\beta_k=\frac{\beta}{(k+K_1)^{\bfrak}},$$
where $0.5<\afrak<\bfrak<1$, $\alpha,\beta\leq 0.5$ and $K_1\geq 1$. This ensures that $(\alpha_k)$ is non-summable $(\sum_k\alpha_k=\infty)$, square-summable $(\sum_k\alpha^2_k<\infty)$ and non-increasing $(\alpha_{k+1}\leq \alpha_k)$, and analogously for $(\beta_k)$. This also ensures that $\lim_{k\uparrow\infty} \beta_k/\alpha_k=0$. Finally, we make the key assumption that 
$\beta_k^2/\alpha_k^3\leq 1.$
\end{assumption}

\begin{lemma}\label{lemma:wp1_and_bound}
    Suppose the setting of Theorem \ref{thm:main} holds. Then,
    \begin{enumerate}[label=\alph*)]
        \item $\|x_k-\xstar(y_k)\|$ converges to zero, and $y_k$ converges to the set $\Ystar$ w.p.\ 1.
        \item There exist constants $C_1,\Gamma_4>0$ such that, for all $k\geq 0$,
        $$\EE[\|x_k-\xstar(y_k)\|^2]\leq \frac{C_1}{(k+1)^{\afrak}}\;\;\text{and}\;\;\sum_{i=0}^k\beta_i\EE\left[\|h(y_i)-y_i\|^2\right]\leq \Gamma_4,$$
        where $C_1=\kappa_2\left(1+\|x_0-\xstar(y_0)\|^2+\|y_0-\ystar\|^2\right)$ for some $\kappa_2>0$.
    \end{enumerate}
\end{lemma}

\begin{align}\label{bound-U}
    \EE[\|y_k-z_k\|^2]= \EE\left[\|U_k\|^2\right]&\leq \sum_{i=0}^{k-1} \left(\beta_i\prod_{j=i+1}^{k-1}(1-\beta_j)\right)^2\EE\left[\|M'_{i+1}\|^2\right]\nonumber\\
    &\stackrel{(a)}{\leq} \sum_{i=0}^{k-1} \beta_i^2\prod_{j=i+1}^{k-1}(1-\beta_j) D_4\nonumber\\
    &\stackrel{(b)}{\leq} 2D_4\beta_k\eqqcolon \Gamma_5\beta_k.
\end{align}

\begin{equation}\label{bound_h-z_i-1}
    \sum_{i=0}^k\beta_i\EE\left[\|h(z_i)-z_i\|^2\right]\leq \Gamma_6.
\end{equation}

\begin{theorem}\label{thm:main}
    Suppose that Assumptions \ref{f-contrac}-\ref{assu:stepsize} hold. Let $\{x_k,y_k\}$ be generated by \eqref{iter-main} with $\beta/\alpha\leq \gamma_1$ and $K_1\geq \gamma_2$. Then, 
    \begin{itemize}
        \item  There exist constants $C_1,C_2>0$ such that for all $k\geq 0$,
        $$\EE\left[\|x_k-\xstar(y_k)\|^2\right]\leq \frac{C_1}{(k+1)^\afrak} \;\;\text{and}\;\; \EE\left[\|g(\xstar(y_k),y_k)-y_k\|^2\right]\leq \frac{C_2}{(k+1)^{1-\bfrak}}.$$
    \item Moreover, $\|x_k-\xstar(y_k)\|$ converges to zero, and the iterates $y_k$ converge to the set $\Ystar$ with probability $1$.
    \end{itemize}
\end{theorem}

\begin{proof}[\textbf{Proof for Theorem \ref{thm:main}}] We have already shown the almost sure convergence and the bound for $\EE[\|x_k-\xstar(y_k)\|^2]$ in Lemma \ref{lemma:wp1_and_bound}. We just need to complete the bound for $\EE[\|h(y_i)-y_i\|^2]$ to complete the proof.
Summing above bound and applying bound from \eqref{bound_h-z_i-1}, we get
\begin{equation}\label{bound_h-z_i-2}
    \EE\left[\|h(z_k)-z_k\|^2\right]\sum_{i=0}^k\beta_i\leq \Gamma_6+\Gamma_7\sum_{i=0}^k\beta_i\sum_{j=i}^{k-1}\beta_j\alpha_j.
\end{equation}

To bound the second term, we first note that 
\begin{align*}
    \sum_{j=i}^{k-1}\beta_j\alpha_j&\leq\sum_{j=i}^{\infty}\frac{\alpha\beta}{(j+K_1)^{\afrak+\bfrak}}\leq\sum_{j=i}^{\infty}\frac{\alpha\beta}{(j+1)^{\afrak+\bfrak}}\leq\alpha\beta\left(\frac{1}{(i+1)^{\afrak+\bfrak}}+\int_{x=i+1}^\infty \frac{1}{x^{\afrak+\bfrak}}dx\right)\\
    &\stackrel{(a)}{\leq} \alpha\beta\left(\frac{1}{(i+1)^{\afrak+\bfrak}}+\frac{1}{(\afrak+\bfrak-1)(i+1)^{\afrak+\bfrak-1}}\right)\stackrel{(b)}{\leq} \frac{2\alpha\beta}{\afrak+\bfrak-1}\frac{1}{(i+1)^{\afrak+\bfrak-1}}.
\end{align*}
Here, inequality (a) and (b) both follow from the fact that $1< \afrak+\bfrak<2$, which follows from our assumption that $\alpha_k$ and $\beta_k$ are both non-summable but square summable sequences. Next, 
\begin{align*}
    &\sum_{i=0}^k\beta_i\sum_{j=i}^{k-1}\beta_j\alpha_j\leq \frac{2\alpha\beta^2}{\afrak+\bfrak-1}\sum_{i=0}^{k-1}\frac{1}{(i+1)^{\afrak+2\bfrak-1}}\stackrel{(c)}{\leq} \frac{2\alpha\beta^2}{\afrak+\bfrak-1}\sum_{i=0}^\infty\frac{1}{(i+1)^{4\afrak-1}}\stackrel{(d)}{\eqqcolon} D_5.
\end{align*}
 Inequality (c) follows from the fact that $2\bfrak\geq 3\afrak$, which follows from our assumption that $\beta_k^2\leq \alpha_k^3$. Finally, inequality (d) follows from the fact that $4\afrak-1>1$, which follows from our assumption that $\afrak>0.5$ or $\alpha_k$ is square summable. 

Hence the right hand side of the inequality \eqref{bound_h-z_i-2} is bounded by a constant. Note that $\sum_{i=0}^k \beta_i\geq (k+1)\beta_k,$ which follows from our assumption that the sequence $\beta_k$ is non-increasing. This implies that 
$$\EE\left[\|h(z_k)-z_k\|^2\right]\leq \frac{\Gamma_6+\Gamma_7D_5}{(k+1)\beta_k}.$$ 
To complete our proof, we recall that $\|h(y_k)-y_k\|^2\leq 2\|h(z_k)-z_k\|^2+8\|y_k-z_k\|^2.$
This implies that 
\begin{align}\label{final-bound-y}
   \EE\left[\|h(y_k)-y_k\|^2\right] &\leq \frac{2\Gamma_6+2\Gamma_7D_5}{(k+1)\beta_k}+8\Gamma_5\beta_k\leq \frac{C_2}{(k+1)^{1-\bfrak}},
\end{align}
where $C_2\coloneqq (2\Gamma_6+2\Gamma_7D_5)/\beta+8\Gamma_5\beta$. Here, the first inequality follows from \eqref{bound-U}, and the last inequality follows from the fact that $1/(k+1)^{\bfrak}\leq 1/(k+1)^{1-\bfrak}$, which follows from our assumption that $\bfrak>1/2$. This completes the proof.
\end{proof}

\end{document}
